{\centering \section*{Lời nói đầu}}













Qua thời gian học tập và rèn luyện tại Đại học Bách khoa Hà Nội, 
đến nay em đã hoàn thành đồ án tốt nghiệp. 
Trong thời gian học tập để có được kết quả hiện tại 
em xin chân thành cảm ơn tới các thầy cô giảng viên trong \textit{Khoa Toán Tin}
đã giảng dạy, quan tâm và tạo điều kiện thuận lợi để em học tập rèn luyện trong suốt thời gian qua, 
giúp em trang bị những kiến thức, kỹ năng cần thiết 
để em có thể hoàn thành tốt nhất báo cáo này 
và phục vụ cho công việc thực tế sau này.

Đồng thời, em xin gửi lời cảm ơn chân thành và sâu sắc đến    \textit{Thầy TS. Trần Anh Tú}, 
người thầy đã tận tình hỗ trợ, 
giúp đỡ và hướng dẫn em suốt thời gian thực hiện đồ án. 
Những kiến thức và kinh nghiệm mà em đã tiếp thu được trong quá trình này 
sẽ đóng góp quan trọng vào sự phát triển và thành công của em trong tương lai. 
Em xin gửi lời chúc sức khỏe tốt nhất đến thầy, 
hy vọng thầy luôn dồi dào sức khỏe, đam mê và nhiệt huyết trong công việc giảng dạy.




Em cũng xin gửi lời cảm ơn tới bố mẹ em 
đã dành nhiều thời gian chăm lo, động viên 
và giúp đỡ em để em đạt được thành công trong học tập.

Mặc dù đã cố gắng nỗ lực học hỏi để hoàn thành đồ án, 
vậy nhưng thời gian thực hiện đồ án có hạn, 
kiến thức của em còn hạn chế. 
Do vậy, 
không tránh khỏi những thiếu sót, 
em rất mong nhận được những ý kiến đóng góp quý báu của thầy cô 
để em có thể cải thiện một cách tốt nhất.


 









\vspace{0.5cm} \emph{Em xin chân thành cảm ơn!}





 

 

Hà Nội, ngày ... tháng ... năm 2026

Sinh viên thực hiện

Vũ Văn Nghĩa




%%%%%%%%%%%%%%%%%%%%%%%%%%%%%%%%%%%%%%%%%%%%%%%%%%%%%%%
%%%%%%%%%%%%%%%%%%%%%%%%%%%%%%%%%%%%%%%%%%%%%%%%%%%%%%%
%%%%%%%%%%%%%%%%%%%%%%%%%%%%%%%%%%%%%%%%%%%%%%%%%%%%%%%

{\centering \section*{Tóm tắt nội dung đồ án}}







Tóm tắt nội dung của đồ án tốt nghiệp trong khoảng tối đa 300 chữ. 
Phần tóm tắt cần nêu được các ý: 
vấn đề cần thực hiện; 
phương pháp thực hiện; 
công cụ sử dụng (phần mềm, phần cứng…); 
kết quả của đồ án có phù hợp với các vấn đề đã đặt ra hay không; 
tính thực tế của đồ án, 
định hướng phát triển mở rộng của đồ án (nếu có); 
các kiến thức và kỹ năng mà sinh viên đã đạt được.








Các chương trong báo cáo được trình bày theo thứ tự sau:

\vspace{0.5cm}

\begin{itemize}
\item
Giới thiệu về công ty

Báo cáo này bắt đầu với phần
giới thiệu về công ty,
bao gồm:
lịch sử hình thành,
sứ mệnh,
slogan,
linh vật biểu tượng,
giá trị cốt lõi,
các sản phẩm nổi bật, \dots

\item
Dự án bài đánh giá kỹ năng Backend ban đầu

Chương này đề cập đến
các kiến thức và kỹ năng
có liên quan và áp dụng cho dự án thực tế.
Đặc biệt là ứng dụng mật mã TOTP xác thực đa yếu tố cho hệ thống.

\item
Dự án DNS Sinkhole Demo

Chương này đề cập đến
các kiến thức về DNS
và kiến thức về DNS Sinkhole.
Từ đó,
áp dụng cho việc chặn quảng cáo,
chặn tên miền không mong muốn,
chặn tên miền độc hại.

\item
Dự án Kong API Gateway

Chương này đề cập đến
các kiến thức về API Gateway
và kiến thức về Kong API Gateway.
Từ đó,
sử dụng Kong
xây dựng
API Gateway
cho
hệ thống microservices.

\item
Dự án thực tế PXDR

Chương cuối cùng này đề cập đến
dự án thực tế
mà em được tham gia thực tập.
Một số công việc nổi bật
em đã học tập, tìm hiểu
và tham gia phát triển
trong dự án.

\end{itemize}

\begin{flushright}
\begin{tabular}{c}
Hà Nội, \today \\
Sinh viên thực hiện \\
\textit{(Ký và ghi rõ họ tên)} \\
Nghĩa \\
Vũ Văn Nghĩa
\end{tabular}
\end{flushright}